\subsection{Density Matrix of the Kitaev Ground State on a Torus}\label{sec:density-matrix-calculate}

So far, we have known two crucial points for constructing the density matrix. The first one is that finite spin-correlation functions of the Kitaev ground state must be of the following form
\[
\anglelr{\prod_{\anglelr{ij}\in \mathfrak{D}}\sigma^{\alpha_{ij}}_i\sigma^{\alpha_{ij}}_j} = \anglelr{\Sigma_{\mathfrak{D}}} \equiv \mathfrak{C}_{\mathfrak{D}}
\]
Otherwise, the correlation function vanishes. From now on, we denote by $\mathfrak{C}_{\mathfrak{D}}$ the expectation values of the string operator corresponding to the edge configuration $\mathfrak{D}$. The other is that every edge configuration in the same equivalence class gives the same weight to the density matrix, so we can sort terms in the expansion \eqref{equ.rho-starting} into equivalence classes
\begin{equation}\label{equ.class-sum}
    \sum_{\mathfrak{D}}\mathfrak{C}_{\mathfrak{D}}\Sigma_\mathfrak{D} = \sum_{[\mathfrak{D}^c]}\sum_{\mathfrak{D}\in[\mathfrak{D}^c]}\mathfrak{C}_{\mathfrak{D}}\Sigma_\mathfrak{D}
\end{equation}
From \eqref{equ.string-operator-equivalent}, we also know that every term of \eqref{equ.class-sum} in an equivalence class can be related to each other via the projector $\mathcal{P}^{(p)}$. Altogether, applying a projector $\mathcal{P}^{(p)}$ is exactly factorizing the density matrix into groups that include every edge configuration in an equivalence class $\sum_{\mathfrak{D}\in [\mathfrak{D}^c]}$, and the form of the density matrix should contain only characters of different classes $[\mathfrak{D}^c]$. Therefore, we anticipate that the form of the density matrix is of the following form
\begin{equation}\label{equ.rho-factorization}
    \rho = C \mathcal{P}_{\mrm{topo}} \sum_{[\mathfrak{D}]} \mathcal{P}^{(p)} \mathfrak{C}_{\mathfrak{D}} \Sigma_{\mathfrak{D}}
\end{equation}
where the topological Wilson-loop projector 
\[
\mathcal{P}_{\mrm{topo}}\equiv \lr{\frac{\mathds{1}+\Tilde{W}_h}{2}}\lr{\frac{\mathds{1}+\Tilde{W}_v}{2}}
\]
is imposed to include configurations related by global loops $\mathcal{L}_v$ and $\mathcal{L}_h$ illustrated in \autoref{fig.equivalent}. The constant $C$ is the normalization factor, which will be determined by imposing the constraint $\Tr\rho = 1$. Carrying out the trace 

\begin{align}\label{equ.rho-direct-trace-1}
    \Tr\rho &= \frac{C}{4} \sum_{[\mathfrak{D}]} \mathfrak{C}_{\mathfrak{D}} \Tr\lr{\mathcal{P}^{(p)}\Sigma_{\mathfrak{D}}} \nonumber\\
    &= \frac{C}{2^{2+N_p}} \sum_{[\mathfrak{D}]} \sum_{\mathcal{A}}'\Tr\lrm{ \lr{2\prod_{p\in\mathcal{A}} W_p}\Sigma_{\mathfrak{D}}}
\end{align}
The coefficient $1/4$ in the first line originates from the denominator of $\mathcal{P}_{\mrm{topo}}$. The two global Wilson-loop operators drop out because neither the local plaquette operators nor the string operators can annihilate them to yield the identity; consequently, their contributions vanish upon taking the trace. The summation of plaquette operators arises from the expansion of $\mathcal{P}^{(p)}$
\begin{align}\label{equ.W-prod-expansion}
    \prod_{p}\lr{\mathds{1}+W_p} &= \lr{\mathds{1}+\prod_{p}W_p}\sum'_{\mathcal{A}}\prod_{p\in\mathcal{A}}W_p \nonumber\\
    &= 2\sum'_{\mathcal{A}}\prod_{p\in\mathcal{A}}W_p.
\end{align}
with $\mathcal{A}$ denoting different sets of plaqette configurations defined in \eqref{equ.lattice-Stokes}. We point out that, due to the identity $\prod_{p}W_p=\mathds{1}$, the loop constructed by $\prod_{p\in \mathcal{A}}W_p$ is identical to that constructed by $\prod_{p\in \overline{\mathcal{A}}}W_p$ with $\overline{\mathcal{A}}$ representing the complement of $\mathcal{A}$. Therefore, we only need to sum over distinct loops, and the sum with a prime symbol on top in \eqref{equ.W-prod-expansion} precisely stands for such a summation. The only two terms that contribute to the trace in \eqref{equ.rho-direct-trace-1} are (i) $\mathcal{A}=\varnothing$ and $\mathfrak{D}\sim \varnothing$ and (ii) $\mathcal{A}=\mathfrak{P}$ and $\mathfrak{D}\sim\mathfrak{D}_Z$. Therefore, \eqref{equ.rho-direct-trace-1} is equal to 
\[
\Tr\rho = \frac{C}{2^{N_p}}\Tr\mathds{1} = \frac{C}{2^{-N+N_p}}
\]
To make the trace equal to $1$, we require $C = 2^{-N+N_p}$. In the end, we obtain the exact form of the density matrix
\begin{equation}\label{equ.density-matrix}
    \rho =  \frac{1}{2^{N-N_p}}\mathcal{P}_{\mrm{topo}}\mathcal{P}^{(p)}\sum_{[\mathfrak{D}]}\mathfrak{C}_{\mathfrak{D}}\Sigma_{\mathfrak{D}}
\end{equation}
Because $[W_p,\Sigma_\mathfrak{D}]=0\;\forall p,\mathfrak{D}$, the projectors $\mathcal{P}^{(p)}$ and $\mathcal{P}_{\mrm{topo}}$ commute with the summation part consisting of string operators. This is also a manifestation of the separability of the $c$-Majorana fermion sector from the gauge sector, which will be elucidated in later sections.

